\textbf{P5)} Dos péndulos idénticos de longitud $\ell$ y masa $m$ oscilan en un mismo plano. Sus masas están unidas por un resorte de constante elástica $k$, el cual se encuentra sin deformar cuando ambos péndulos están en su posición de equilibrio vertical.

Se considera que las oscilaciones son pequeñas, de modo que se puede usar la aproximación $\sin\theta\approx\theta$.

Sea $x_1(t)$ y $x_2(t)$ los desplazamientos horizontales de cada masa respecto a su posición de equilibrio.

\begin{enumerate}[a)]
  \item Usando la aproximación de pequeñas oscilaciones, derive las ecuaciones de movimiento para cada masa.
  \item Muestre que, en ausencia del resorte ($k=0$), cada péndulo se comporta como un movimiento armónico simple e identifique su frecuencia angular natural
  \[\omega_0=\sqrt{\frac{g}{\ell}}.\]
  \item Considere ahora el efecto del resorte. Explique cómo la interacción entre ambos péndulos modifica el movimiento respecto al caso de un MAS simple.
  \item Para pequeñas oscilaciones, proponga soluciones del tipo
  \[x_i(t)=A_i\cos(\omega t)\]
  y determine las posibles frecuencias normales del sistema (modo simétrico y modo antisimétrico).
  \item Describa físicamente los distintos tipos de movimiento que pueden ocurrir (por ejemplo, ambos péndulos en fase —modo simétrico— o en oposición de fase —modo antisimétrico—) y cómo se relacionan con las frecuencias encontradas.
\end{enumerate}

\vspace{0.3cm}
